\section{Introduction}
\label{sec:intro}

% Paragraph 1 -- the setting. What problem does the field care about, and what
% is the current state of practice? Two or three sentences, no hedging.
\ARTODO{setting and why it matters}

% Paragraph 2 -- the gap. Be specific about what existing work does *not* do.
% This paragraph is where a reviewer decides whether the paper is novel, so it
% must name concrete prior approaches and their concrete limitation.
\ARTODO{the specific gap, naming the prior approaches it belongs to}

% Paragraph 3 -- the idea, in plain language before any notation.
\ARTODO{one-paragraph statement of the idea}

% Paragraph 4 -- what the experiments show. Every number here must trace back
% to a row in Section~\ref{sec:experiments}; leave \ARnum standing until it does.
\ARTODO{headline empirical claim, with \ARnum{} where the number will go}

% Paragraph 5 -- contributions. Three or four items, each one a claim the paper
% actually supports. Reviewers check these against the results table.
\paragraph{Contributions.}
\begin{itemize}
  \item \ARTODO{contribution 1: the method or formulation}
  \item \ARTODO{contribution 2: the analysis or theory, if any}
  \item \ARTODO{contribution 3: the empirical result, with the metric named}
\end{itemize}
